%
%      Brno University of Technology
%      Faculty of Information Technology
%
%      MSc Thesis	2009/2010
%
%      Design of S-Boxes Using Genetic Algorithms
%
%
%      Bedrich Hovorka
%
% -----------------------------------------------------
Cílem experimentů bylo zjistit zda jsou Evoluční algoritmy
vhodné pro pro hledání zvolených velikostí s-boxů.
Za tím účelem bylo porovnáno s náhodným prohledáváním.
Porovnání běhů dle statistik z Galibu \cite{galib} na několika bězích.
Experimenty byly naskriptovány pomocí Pythonu.
Byly vybrány na základě testovaní programu, kde proběhly jednoběhové zkoušky
velikosti populace. Pro dlouhodobější hledání kvalitnějších řešení byly nejprve
stanoveny parametry evoluce, zejména pravděpodobnosti mutace a křížení.
Poté se tyto parametry užily pro náročnější experimenty.

\section{Stanovení pravděpodobností mutace a křížení pro jednotlivá kritéria}
Nejprve byly stanoveny parametry křížení a mutace pro jednokriteriální
optimalizaci pro každé kritérium. Při nich se většinou dospělo k mezním hodnotám.
Byly nalezeny nejlepší pravděpodobnosti mutace a křížení pro každé kritérium
při pevné velikosti populace 100, počtu generací 1000 a počtu běhů 100.
Proti tomu bylo postaveno stejně početné náhodné prohledávání.
S-box byl zvolen středně složitý se 4 vstupy a 4 výstupy (tedy byl typu $4 \times 4$).
Nejprve byly na intervalu $\langle 0, 1)$ zkoušeny všechny možné dvojice pravděpodobností:
pro mutaci $0, 0.05, 0.1, 0.2, 0.4, 0.6, 0.8$; pro křížení $0, 0.2, 0.4, 0.6, 0.8$.
Protože výsledky u některých kritérií ve vyšších hodnotách byly obdobné a málo slibnější,
byl zvolen ještě užší interval $\langle0, 0.2)$. Zkoušeny byly
dvojice pravděpodobností $0, 0.04, 0.08, 0.12, 0.16$ jak u křížení tak u mutace.
U každého běhu bylo zaznamenáno uloženo 100 nejlepších jedinců z průběhu celé evoluce
a všichni jedinci ze všech běhů pro jednu dvojici $(p_{Mut}, p_{Cross})$ byly sloučeni
dohromady a nad nimy byly spočteny statistiky.
Výsledky zaneseny do tabulek, které jsou seřazeny od nejhoršího po nejlepší
podle řadícího klíče (maximum, medián, průměr + odchylka, průměr, minimum).
Pro rozhodnutí byl proveden orientační dvouvýběrový t-test výsledků náhodného prohledávání
a nejlepší kombinace parametrů pro genetický algoritmus.
\subsubsection{Bent skóre}
\begin{table}[tb]
\begin{center}
\begin{tabular}{| l | c | c | c | c | c |}
\hline
Parametry & maximum & medián & std. odchylka & průměr & minimum \\ \hline
(0, 0) & $32$ & $31$ & $4.07067$ & $29.1836$ & $7$ \\ \hline
ostatní kombinace & $32$ & $32$ & $0$ & $32$ & $32$ \\ \hline
\end{tabular}
\caption{Výsledky pro celý interval a kritérium bent skóre $+$ řád SAC}
\label{tabBENT_AND_MOSACParameters}
\end{center}
\end{table}
V tabulce \ref{tabBENT_AND_MOSACParameters} je uvedena pouze varianta bez
křížení a mutace, ostatní kombinace parametrů dávaly stejné statistiky
a to takové, že se shodovaly i se statistikami náhodného prohledávání,
u kterého vyšlo maximum 32, medián 32, odchylka 0, průměr 32, minimum 32.
Tj. všichni nejlepší členové dosáhli maximálního možného bent skóre pro
bijektivní s-box. Pro bijektivní by to bylo 64.
U tohoto kritéria tedy nezáleží na parametrech pro velikost
populace 100 a počet generací 1000. Výsledky se shodují s náhodným prohledáváním.
Na základě toho se kritérium dále pro zúžení intervalu zdálo nezajímavé.

\subsubsection{Linear probability}
\begin{table}[tb]
\begin{center}
\begin{tabular}{| l | c | c | c | c | c |}
\hline
Parametry & maximum & medián & std. odchylka & průměr & minimum \\ \hline
(0, 0) & $2$ & $0.830075$ & $0.380956$ & $0.910228$ & $0$ \\ \hline
ostatní kombinace & $2$ & $2$ & $0$ & $2$ & $2$ \\ \hline
(0, 0.2) & $2$ & $2$ & $0.452339$ & $1.7859$ & $0.830075$ \\ \hline
\end{tabular}
\caption{Výsledky pro celý interval pro kritérium $LP_{\max}$}
\label{tabLP_MAXParameters}
\end{center}
\end{table}
\begin{table}[tb]
\begin{center}
\begin{tabular}{| l | c | c | c | c | c |}
\hline
Parametry & maximum & medián & std. odchylka & průměr & minimum \\ \hline
(0, 0) & $2$ & $0.830075$ & $0.373374$ & $0.905579$ & $0$ \\ \hline
(0, 0.04) & $2$ & $0.830075$ & $0.511175$ & $1.13063$ & $0.830075$ \\ \hline
(0, 0.08) & $2$ & $0.830075$ & $0.577091$ & $1.31957$ & $0.830075$ \\ \hline
ostatní kombinace & $2$ & $2$ & $0$ & $2$ & $2$ \\ \hline
(0, 0.12) & $2$ & $2$ & $0.584362$ & $1.44218$ & $0.830075$ \\ \hline
(0, 0.16) & $2$ & $2$ & $0.549793$ & $1.61486$ & $0.830075$ \\ \hline
\end{tabular}
\caption{Výsledky pro zúžený interval pro kritérium $LP_{\max}$} %LP_MAX
\label{tabLP_MAXParametersSmall}
\end{center}
\end{table}

V tabulce \ref{tabLP_MAXParameters} vychází jako nejvyšší hodnota $2$.
Maximální odchylka pravděpodobnosti lineárních výrazů $2^{-2} = 0.25$.
Vyzkoušeny hodnoty v užším intervalu a zaznamenány do tabulky \ref{tabLP_MAXParametersSmall}.
Náhodné prohledávání mělo: maximum 2, medián 2, odchylka 0.365039, průměr 1.87213, minimum 0.830075.
Náhodné prohledávání má lepší průměr.

\subsubsection{Diferential probability}
\begin{table}[tb]
\begin{center}
\begin{tabular}{| l | c | c | c | c | c |}
\hline
Parametry & maximum & medián & std. odchylka & průměr & minimum \\ \hline
(0, 0) & $2$ & $1.41504$ & $0.280306$ & $1.28179$ & $0.415038$ \\ \hline
ostatní kombinace & $2$ & $2$ & $0$ & $2$ & $2$ \\ \hline
(0, 0.2) & $2$ & $2$ & $0.291124$ & $1.73563$ & $1.41504$ \\ \hline
\end{tabular}
\caption{Výsledky pro celý interval pro kritérium $DP_{\max}$}
\label{tabDP_MAXParameters}
\end{center}
\end{table}
\begin{table}[tb]
\begin{center}
\begin{tabular}{| l | c | c | c | c | c |}
\hline
Parametry & maximum & medián & std. odchylka & průměr & minimum \\ \hline
(0, 0) & $2$ & $1.41504$ & $0.283811$ & $1.28443$ & $0.415038$ \\ \hline
(0, 0.04) & $2$ & $1.41504$ & $0.199666$ & $1.48975$ & $1$ \\ \hline
(0, 0.08) & $2$ & $1.41504$ & $0.23421$ & $1.53232$ & $1.41504$ \\ \hline
(0, 0.12) & $2$ & $1.41504$ & $0.264779$ & $1.58327$ & $1.41504$ \\ \hline
(0, 0.16) & $2$ & $1.41504$ & $0.284352$ & $1.63902$ & $1.41504$ \\ \hline
ostatní kombinace & $2$ & $2$ & $0$ & $2$ & $2$ \\ \hline
\end{tabular}
\caption{Výsledky pro zúžený interval pro kritérium $DP_{\max}$} %DP_MAX
\label{tabDP_MAXParametersSmall}
\end{center}
\end{table}
Výsledky pro celý interval v tabulce \ref{tabDP_MAXParameters} jsou obdobné výsledkům kritéria $LP_{\max}$.
Náhodné prohledávání mělo: maximum 2, medián 2, odchylku 0.29122, průměr 1.73475 a minimum 1.41504.
Pro zúžený interval tabulka \ref{tabDP_MAXParametersSmall} potvrzuje, že pokud je vypnutá mutace
a křížení je menší než 0.2 nalezne více slabších výsledků. V jiném případě je lepší než náhodné prohledávání.

\pagebreak
\subsubsection{SAC} %SAC
Samotný řád SAC má malou vypovídací hodnotu o kvalitě řešení, přesto
genetický algoritmus nalezl nejvyšší možný řád. Hodnoty v tabulce \ref{tabSACParameters} jsou zvýšené
o jedničku, aby fitness nebyla záporná\footnote{Což vyžaduje Galib}.
Z tabulky vyplývá, že takřka nezáleží na pravděpodobnosti
mutace a křížení. Rozdíly mezi jednotlivými dvojicemi parametrů jsou nevýznamné.
Naproti tomu náhodné prohledávání dospělo k maximu 1, mediánu 0, odchylce 0.249148,
průměru 0.0665 a minimu 0. Tj. našlo jen pár jedinců, kteří vůbec splňují SAC 0. řádu.

\subsubsection{Branching faktor} %BF
Genetický algoritmus s faktorem větvení jako objektivní funkcí v tabulce \ref{tabBFParameters}
dospělo k obdobným statistikám jako náhodné prohledávání, které
dospělo k maximu 3, mediánu 2, odchylce 0.0199968, průměru 2.0004 a minimu 2.
Pro zúžení byl nezajímavý.

\subsubsection{Stupeň algebraického polynomu} %POLYNOMIAL_DEGREE
\begin{table}[tb]
\begin{center}
\begin{tabular}{| l | c | c | c | c | c |}
\hline
Parametry & maximum & medián & std. odchylka & průměr & minimum \\ \hline
ostatní kombinace & $3$ & $3$ & $0$ & $3$ & $3$ \\ \hline
(0, 0) & $3$ & $3$ & $0.457029$ & $2.748$ & $1$ \\ \hline
\end{tabular}
\caption{Výsledky pro celý interval pro kritérium stupeň algebraického polynomu} %POLYNOMIAL_DEGREE
\label{tabPOLYNOMIAL_DEGREEParameters}
\end{center}
\end{table}
\begin{table}[bt]
\begin{center}
\begin{tabular}{| l | c | c | c | c | c |}
\hline
Parametry & maximum & medián & std. odchylka & průměr & minimum \\ \hline
ostatní kombinace & $3$ & $3$ & $0$ & $3$ & $3$ \\ \hline
(0, 0) & $3$ & $3$ & $0.455533$ & $2.7484$ & $1$ \\ \hline
\end{tabular}
\caption{Výsledky pro zúžený interval pro kritérium stupeň algebraického polynomu} %POLYNOMIAL_DEGREE
\label{tabPOLYNOMIAL_DEGREEParametersSmall}
\end{center}
\end{table}
Všechny kombinace parametrů v tabulce \ref{tabPOLYNOMIAL_DEGREEParameters} až na variantu
bez křížení a mutace se ve výsledcích shodují.
Stálo za to, prověřit to i v zúženém intervalu v tabulce \ref{tabPOLYNOMIAL_DEGREEParametersSmall}, který
zjištění potvrzuje.
Náhodné prohledávání dospělo též maximu 3, mediánu 3, odchylce 0, průměru 3 a minimu 3.
Pro bijektivní funkce je to stupeň polynomu nejvyšší možný.
\pagebreak
\begin{table}[bt]
\begin{center}
\begin{tabular}{| l | c | c | c |}
\hline
Kritérium & pMut & pCross & co je lepší \\ \hline
bent skóre & 0.1 & 0.2 & Random \\ \hline
$LP_{\max}$ & 0.0 & 0.16 & Random \\ \hline
$DP_{\max}$ & 0.0 & 0.2  & GA \\ \hline
SAC & 0.1 & 0.2 & GA \\ \hline
faktor větvení & 0.1 & 0.2 & Random \\ \hline
stupeň polynomu & 0 & 0 & Random \\ \hline
\end{tabular}
\caption{Shrnutí Genetického algoritmu vs. náhodného prohledávání na permutaci}
\label{tabGAvsRandom}
\end{center}
\end{table}

V tabulce \ref{tabGAvsRandom} jsou uvedeny nejlepší možné parametry genetického algoritmu.
A stanovení, který algoritmus byl pro provedené evoluce lepší.
Tam kde se hodnoty statistik shodovaly bylo zvoleno náhodné prohledávání kvůli své jednodušší
algoritmické složitosti. Na statistikách evoluce různých jsou vidět jak omezení kritérií, tak
i podobný tvar vzorců, kterými jsou k kritéria počítána.
To jaký je algoritmus lepší záleží na použitém kritériu.
V některých případech má evoluci smysl provést.

\section{Porovnání výsledků EDA algoritmů\\s Genetickým algoritmem}
Dalším provedeným experimentem bylo ověření na binárním chromozómu,
zda EDA algoritmy nacházejí lepší, horší či stejné výsledky v porovnání s nejlepšími parametry předchozího
experimentu. Byly provedeny stejné výpočty jako v předešlém experimentu, tentokrát však pro
nejlepší parametry Genetického algoritmu vůči algoritmům BMDA a UMDA a ná\-hod\-né\-mu prohledávání.
Narozdíl od předchozího však nebyla hledána bijektivní funkce.

Pro vizualizaci porovnání byly zvoleny tabulky kratší, které ukazují
kritéria v jednotlivých algoritmech, a krabicové grafy, které porovnávají
algoritmy mezi sebou. Krabicový graf znázorňuje rozložení stejně početných skupin hodnot.
Horní a dolní strana modrého obdélníka znázorňuje umístění kvartilů. Prostřední červená
vodorovná úsečka přestavuje medián. Nahoru a dolů od obdélníka vede tzv. vous \footnote{z anglického whisker}.
Horní resp. dolní vous znázorňuje nejvyšší resp. nejnižší hodnotu, na kterou má být brán zřetel.
Křížky jsou označeny odlehlé hodnoty.
V tabulce \ref{tabGaGaVsEDAComparison} jsou statistiky pro Genetický algoritmus.
Pro SAC se nepodařilo najít nejvyšší řád. Pro Bent skóre a stupeň polynomu již neplatí
omezení bijektivitou. Jinak jsou výsledky podobné hledání permutace z předcházejícího
experimentu. UMDA algoritmus v tabulce \ref{tabUMDAGaVsEDAComparison}
neměl žádný problém z $DP_{\max}$ a stupněm polynomu, za to SAC podobně jako u náhodného prohledávání
našel jen pár jedinců, kteří splňovali nejnižší stupeň SAC.
BDMA v tabulce \ref{tabBMDAGaVsEDAComparison} dopadl obdobně, trochu lepší výsledek měl
u bent skóre. Náhodné prohledávání (tab. \ref{tabRandomGaVsEDAComparison}) mělo stejné výsledky v  maximech a mediánech jako
BMDA.%az jsem pojal podezreni ze jsem to blbe zkopiroval z vystupu programu

\begin{table}[tb]
\begin{center}
\begin{tabular}{| l | c | c | c | c | c |}
\hline
Kritérium & maximum & medián & std. odchylka & průměr & minimum \\ \hline
$LP_{\max}$ & $2$ & $1.35614$ & $0.221604$ & $1.44462$ & $1.35614$ \\ \hline
bent skóre & $64$ & $40$ & $3.51839$ & $42.0224$ & $40$ \\ \hline
$DP_{\max}$ & $2$ & $2$ & $0.0628992$ & $1.99316$ & $1.41504$ \\ \hline
SAC & $2$ & $1$ & $0.0141407$ & $1.0002$ & $1$ \\ \hline
BF & $2$ & $2$ & $0$ & $2$ & $2$ \\ \hline
stupeň polynomu & $4$ & $3$ & $0.421644$ & $2.9443$ & $0$ \\ \hline
\end{tabular}
\caption{Porovnání EDA a GA na binárním chromozómu: kritéria pro GA}
\label{tabGaGaVsEDAComparison}
\end{center}
\end{table}
\nopagebreak
\begin{table}[tb]
\begin{center}
\begin{tabular}{| l | c | c | c | c | c |}
\hline
Kritérium & maximum & medián & std. odchylka & průměr & minimum \\ \hline
$LP_{\max}$ & $2$ & $1.35614$ & $0.141776$ & $1.38912$ & $1.35614$ \\ \hline
bent skóre & $47$ & $31$ & $2.23823$ & $31.988$ & $30$ \\ \hline
$DP_{\max}$ & $2$ & $2$ & $0$ & $2$ & $2$ \\ \hline
SAC & $1$ & $0$ & $0.0990082$ & $0.0099$ & $0$ \\ \hline
BF & $2$ & $2$ & $0.498952$ & $1.5322$ & $1$ \\ \hline
stupeň polynomu & $4$ & $4$ & $0$ & $4$ & $4$ \\ \hline
\end{tabular}
\caption{Porovnání EDA a GA na binárním chromozómu: kritéria pro UMDA}
\label{tabUMDAGaVsEDAComparison}
\end{center}
\end{table}
\nopagebreak
\begin{table}[tb]
\begin{center}
\begin{tabular}{| l | c | c | c | c | c |}
\hline
Kritérium & maximum & medián & std. odchylka & průměr & minimum \\ \hline
$LP_{\max}$ & $2$ & $1.35614$ & $0.234303$ & $1.45737$ & $1.35614$ \\ \hline
bent skóre & $48$ & $32$ & $2.98262$ & $33.5576$ & $31$ \\ \hline
$DP_{\max}$ & $2$ & $2$ & $0$ & $2$ & $2$ \\ \hline
SAC & $1$ & $0$ & $0.196198$ & $0.0401$ & $0$ \\ \hline
BF & $2$ & $2$ & $0$ & $2$ & $2$ \\ \hline
stupeň polynomu & $4$ & $4$ & $0$ & $4$ & $4$ \\ \hline
\end{tabular}
\caption{Porovnání EDA a GA na binárním chromozómu: kritéria pro BMDA}
\label{tabBMDAGaVsEDAComparison}
\end{center}
\end{table}
\nopagebreak
\begin{table}[tb]
\begin{center}
\begin{tabular}{| l | c | c | c | c | c |}
\hline
Kritérium & maximum & medián & std. odchylka & průměr & minimum \\ \hline
$LP_{\max}$& $2$ & $1.35614$ & $0.184163$ & $1.4141$ & $1.35614$ \\ \hline
bent skóre & $48$ & $32$ & $4.16932$ & $30.4523$ & $24$ \\ \hline
$DP_{\max}$ & $2$ & $2$ & $0$ & $2$ & $2$ \\ \hline
SAC & $1$ & $0$ & $0.205328$ & $0.0441$ & $0$ \\ \hline
BF & $2$ & $2$ & $0$ & $2$ & $2$ \\ \hline
stupeň polynomu & $4$ & $4$ & $0.478596$ & $3.6447$ & $3$ \\ \hline
\end{tabular}
\caption{Porovnání EDA a GA na binárním chromozómu: kritéria pro Random}
\label{tabRandomGaVsEDAComparison}
\end{center}
\end{table}

\insertimage[0.85\textwidth]{edaVsGaComparisonBENT_AND_MOSAC}{Porovnání EDA algoritmů na binárním reprezentaci s-boxu, kritérium bent skóre} %BENT_AND_MOSAC

Krabicový graf na obrázku \ref{edaVsGaComparisonBENT_AND_MOSAC} porovnává jednotlivé
algoritmy pro kritérium bent skóre. V tomto kritériu je zcela jasným vítězem
Genetický algoritmus. Náhodné prohledávání je v tomto kritériu srovnatelné s EDA algoritmy.
V kritériích $LP_{\max}$ a $DP_{\max}$ (obr. \ref{edaVsGaComparisonLP_MAX} a \ref{edaVsGaComparisonDP_MAX})
nejsou mezi algoritmy rozdíly. U lavinového efektu, kde má genetický algoritmus svůj průměr i medián, jsou
u ostatních odlehlé maximální hodnoty (obr. \ref{edaVsGaComparisonSAC}).
Ve faktoru větvení (obr. \ref{edaVsGaComparisonBF}) je o něco horší UMDA, ale jinak se
dospělo ke stejným výsledkům. Za to ve stupni polynomu (obr. \ref{edaVsGaComparisonPOLYNOMIAL_DEGREE}) je genetický algoritmus
nejhorší. Jen v tomto kritériu vycházejí EDA algoritmy o něco málo lépe než náhodné
prohledávání, leč střední hodnota je shodná.

Celkově tedy EDA algoritmy dopadly nejhůř, náhodné prohledávání nalezne obdobné výsledky.
Původní motivací pro nasazení EDA algoritmů
byly tvary MOSAC Booleovských funkcí, které mají určité
pravidelnosti \footnote{viz tabulka \ref{miniSAC}, větší funkce v \cite{yokohama}}.
Dvouvstupová booleovská funkce (půl-bajt) splňující MOSAC obsahuje 3 stejné bity, čtvrtý se liší.
Třívstupová (bajt) se skládá z dvouvstupových, tak že pozice lišícího se bitu v půl-bajtu jsou symetrické
podle osy vedoucí středem bajtu. U čtyřstup nesmí být dva vedlejší bajty symetrické.
Otázkou by bylo zda-li se z toho dá vytvořit nějaký pravděpodobnostní model závislostí mezi bity,
který by byl použitelný v evoluci pomocí EDA. Proti tomu stojí myšlenka, že
statistický model hledající závislosti mezi bity má pro bezpečnost opačný význam,
právě, že by bity by měly být vzájemně nezávislé, proto ani nebyl implementován
složitější pravděpodobnostní model pro algoritmus BOA.

\section{Parametry paralelního náhodného prohledávání}\label{secParalelRandomComparison}
Pro reprezentaci s-boxu jako acyklický graf (Kartézské genetické programování) není definováno křížení.
Srovnáno bylo s reprezentací i686 u které též nebylo prováděno křížení.
Byly použity dva algoritmy: náhodné prohledávání a paralelní náhodné prohledávání.
Pro paralelní náhodné prohledávání vyzkoušeny byly tyto hodnoty mutace:
0, 0.01, 0.05, 0.1, 0.2, 0.3, 0.4, 0.5, 0.6, 0.7, 0.8 a 0.9.

Chromozóm pro Kartézské genetické programování zvolen se čtyřmi vstupy, čtyřmi výstupy, čtyřmi sloupci,
třemi řádky a s maximálním počtem mutací 6. V tabulce \ref{tabCGPParallelRadnomParameters} jsou shrnuty až výsledné
hodnoty s nejlepšími mutacemi. Hodnoty pro jednotlivá kritéria pro každou hodnotu mutace
jsou zakreslena do grafů na obrázcích \ref{cgpVsLuffaCGPComparisonBENT_AND_MOSAC} až \ref{cgpVsLuffaCGPComparisonPOLYNOMIAL_DEGREE}.
Z nich je patrné, že počet generací je tak velký, že velikost mutace nehraje moc roli.
Pro každé kritérium vycházelo lépe paralelní náhodné prohledávání,
takže tzv. \uv{paralelnost}\footnote{ve smyslu, že je vyhodnoceno více jedinců v populaci}
prokázala svůj smysl.
Pro kritérium $DP_{\max}$ dokonce nalezlo lepší výsledek než předcházející algoritmy.
To samé se dá říci i o faktoru větvení.
\begin{table}[bt]
\begin{center}
\begin{tabular}{| l | c | c | c | c | c | c |}
\hline
Kritérium & pMut & maximum & medián & std. odchylka & průměr & minimum \\ \hline
$LP_{\max}$ & 0.01 & $2$ & $2$ & $0$ & $2$ & $2$ \\ \hline
$DP_{\max}$ & 0.4 & $3$ & $2$ & $0.139994$ & $2.02$ & $2$ \\ \hline
bent skóre & 0.01 & $64$ & $64$ & $0$ & $64$ & $64$ \\ \hline
SAC & 0.01 & $3$ & $3$ & $1.22189$ & $2.37$ & $0$ \\ \hline
BF & 0.9 & $4$ & $3$ & $0.511454$ & $2.72$ & $2$ \\ \hline
stupeň polynomu & 0.01 & $4$ & $4$ & $0$ & $4$ & $4$ \\ \hline
\end{tabular}
\caption{Nejlepší pravděpodobnosti mutace pro paralelní náhodné programování nad reprezentací CGP} %CGP
\label{tabCGPParallelRadnomParameters}
\end{center}
\end{table}
%SOFTWARE_IMPL
\insertimage[0.85\textwidth]{cgpVsLuffaSOFTWARE_IMPLComparisonBENT_AND_MOSAC}{Parametry paralelního náhodného prohledávání,bent skóre, reprezentace Luffa} %BENT_AND_MOSAC
\insertimage[0.85\textwidth]{cgpVsLuffaSOFTWARE_IMPLComparisonDP_MAX}{Parametry paralelního náhodného prohledávání,kritérium DP\_MAX, reprezentace Luffa} %DP_MAX
\begin{table}[bt]
\begin{center}
\begin{tabular}{| l | c | c | c | c | c | c |}
\hline
Kritérium & pMut & maximum & medián & std. odchylka & průměr & minimum \\ \hline
$LP_{\max}$ & 0.5 & $2$ & $2$ & $0.263511$ & $1.86451$ & $0.830075$ \\ \hline
$DP_{\max}$ & 0.05 & $3$ & $2$ & $0.141421$ & $2$ & $1$ \\ \hline
bent skóre & 0.7 & $64$ & $64$ & $4.53159$ & $59.8606$ & $48$ \\ \hline
SAC & 0.05 & $3$ & $3$ & $0.729079$ & $2.78$ & $0$ \\ \hline
BF & 0.05 & $4$ & $4$ & $0.42989$ & $3.7919$ & $2$ \\ \hline
stupeň polynomu & 0.01 & $4$ & $4$ & $0$ & $4$ & $4$ \\ \hline
\end{tabular}
\caption{Nejlepší pravděpodobnosti mutace pro paralelní náhodné programování nad reprezentací i686} %SOFTWARE_IMPL
\label{tabSOFTWARE_IMPLParallelRadnomParameters}
\end{center}
\end{table}

\pagebreak
Pro reprezentaci i686 byly zvolen parametr 1 čtveřice (tj. 1 bit v registru), aby měla
počet vstupů i výstupů rovný čtyřem. Z tabulky \ref{tabSOFTWARE_IMPLParallelRadnomParameters} je patrná
podobnost výsledků těchto reprezentací. Zde však parametr mutace hraje větší roli.
Z krabicového grafu na obrázku \ref{cgpVsLuffaSOFTWARE_IMPLComparisonBENT_AND_MOSAC} je patrné, že
příliš nízká a příliš vysoká pravděpodobnost mutace zhoršuje výsledek pro bent skóre,
hodnoty mezi tím jsou stabilizované. Jako nejlepší parametr však je vybráno 0.7,
protože by při některých jiných kritérií, mohly hodnoty u mutací 0.01 až 0.6 znamenat
uvíznutí v lokálním extrému. To samé platí pro kritérium $DP_{\max}$.
na obrázku \ref{cgpVsLuffaSOFTWARE_IMPLComparisonDP_MAX}. O ostatních kritériích na obrázcích
\ref{cgpVsLuffaSOFTWARE_IMPLComparisonLP_MAX} až \ref{cgpVsLuffaSOFTWARE_IMPLComparisonPOLYNOMIAL_DEGREE}
se dá konstatovat totéž.

Obě reprezentace dospívají k obdobným výsledkům s tím rozdílem, že u každé jsou jiné hodnoty mutace.
To je zřejmě dané strukturou chromozómů a různými způsoby mutování.

\section{Parametry multikriteriálního vyhledávání}
Pro jedno kritérium se tedy dá dospět k praktickým mezním hodnotám rychle a ne moc náročně.
Ovšem jak vypadají výsledky a jaké parametry vyšly nejlépe, když
bude třeba splnit všechna kritéria najednou, probere následující podkapitola.
Nejprve bylo nutné stanovit jak vůbec výsledky pro každou dvojici pravděpodobností mutace a křížení
mezi sebou porovnávat. Nabízel se
počet ve všech nadprůměrných jedinců.
Průměrný jedinec byl vypočten ze všech běhů pro všechny parametry, tj. po provedení
všech potřebných evolucí byly vypočteny průměrné hodnoty pro každé kritérium a z nich
byl vytvořen vektor o stejných dimenzích jako ohodnocení kteréhokoliv skutečného jedince.
Pro každou dvojici pravděpodobností mutace a křížení byli spočteni jedinci,
kteří tomuto \uv{průměrnému jedinci} slabě dominují.
Jako kritéria byla použita všechna, která byla zkoumána v před\-chá\-ze\-jí\-cích podkapitolách,
tedy bent skóre, $LP_{\max}$, $DP_{\max}$, SAC, faktor větvení a stupeň polynomu.
Použitými algoritmy byly VEGA a SPEA.

\subsubsection{Parametry pro permutaci}
Nejprve evoluce proběhla při velikosti populace 100 při 1000 generacích pro permutaci
se čtyřmi vstupy a čtyřmi výstupy.

Algoritmus VEGA dospěl k průměrnému jedinci\footnote{při sázení textu všechny hodnoty zaokrouhleny na dvě desetinná místa}
s bent skóre $29.61$,
$-log_2{LP_{\max}} = 0.99$ tj. $LP_{\max} = 0.5$,
$-log_2{DP_{\max}} = 1.31$ tj. $DP_{\max} = 0.4$,
SAC $0.95$,
faktor větvení $2$
a stupeň polynomu $2.77$.

Algoritmus SPEA dospěl k
bent skóre $31.50$,
$-log_2{LP_{\max}} = 1.91$ tj. $LP_{\max} = 0.27$,
$-log_2{DP_{\max}} = 1.93$ tj. $DP_{\max} = 0.26$,
SAC $0.98$,
faktor větvení $2.04$
a stupeň polynomu $2.96$.

Žádný algoritmus nedospěl k žádnému parametru, co by měl nějaké nadprůměrné jedince.
Proto byly i vzhledem k časové náročnosti algoritmu Spea a na základě výsledků
z testovací fáze programu zvoleny jiné parametry na
počet generací na 200 a velikost populace 200. Počet běhů zůstal na 100.
Hodnoty průměrného jedince byly ještě před stanovením počtů nadprůměrných
zaokrouhleny dolů na celé desetiny.

Algoritmus VEGA dospěl k tomuto průměrnému jedinci s
bent skóre $29.1$,
$-log_2{LP_{\max}} = 1.7$ tj. $LP_{\max} = 0.30778610333622908$,
$-log_2{DP_{\max}} = 1.5$ tj. $DP_{\max} = 0.35355339059327379$,
SAC $0.2$,
faktor větvení $2.0$
a stupeň polynomu $2.7$.
V tabulce \ref{tabVegaParameters} jsou uvedeny výsledky pro jednotlivé kombinace pravděpodobností.
Nejlépe vyšla kombinace s $pMut = 0.3$ a $pCross = 0.8$, počty však byly docela vyrovnané.

Algoritmus SPEA dospěl k o něco lepšímu průměrnému jedinci, podobně jak tomu bylo u vyššího počtu
jedinců v populaci, s
bent skóre $31.6$,
$-log_2{LP_{\max}} = 1.9$ tj. $LP_{\max} = 0.26794336563407328$,
$-log_2{DP_{\max}} = 1.9$ tj. $DP_{\max} = 0.26794336563407328$,
SAC $0.9$,
faktor větvení 2.0
a stupeň polynomu 2.9.
V tabulce \ref{tabSpeaParameters} jsou uvedeny výsledky pro jednotlivé kombinace pravděpodobností.
Nejlépe vyšla kombinace s $pMut = 0.9$ a $pCross = 0.05$.

Zvýšení velikosti populace tedy populace zlepšilo průměry jednotlivých kritérií.

\subsubsection{Parametry pro i686}
Dále proveden předchozí pokus, ale z reprezentací i686 se čtyřmi vstupy
a čtyřmi výstupy, navíc přibylo kritérium bijektivní skóre,
neboť tato reprezentace nezaručuje bijektivitu.

Algoritmus VEGA dospěl k průměrnému jedinci s
bent skóre 14.4,
$-log_2{LP_{\max}} = 0.1$ tj. $LP_{\max} = 0.93303299153680741$,
$-log_2{DP_{\max}} = 0.6$ tj. $DP_{\max} = 0.659753955386447$,
bijektivní skóre 8.8,
SAC 0.0,
faktor větvení 1.0
a stupeň polynomu 1.5.
V tabulce \ref{tabVegaParametersLuffa} jsou uvedeny výsledky pro jednotlivé kombinace pravděpodobností.
Nejlépe vyšla kombinace s $pMut = 0.01$ a $pCross = 0.09$.

Algoritmus SPEA dospěl k průměrnému jedinci s
bent skóre 21.5,
$-log_2{LP_{\max}} = 0.4$ tj. $LP_{\max} = 0.75785828325519899$,
$-log_2{DP_{\max}} = 0.8$ tj. $DP_{\max} =  0.57434917749851744$,
bijektivní skóre 12.2,
SAC 0.4,
faktor větvení 1.5
a stupeň polynomu 2.2.
V tabulce \ref{tabSpeaParametersLuffa} jsou uvedeny výsledky pro jednotlivé kombinace pravděpodobností.
Nejlépe vyšla kombinace s $pMut = 0.05$ a $pCross = 0.7$

Přidání bijektivního skóre má vliv na degradaci ostatních parametrů.

\subsubsection{Parametry pro CGP}
A pokus byl proveden ještě s reprezentací CGP se 4 vstupními bity,
4 výstupními bity, 4 sloupci, 3 řádky a 6 maximálními mutacemi).
Samozřejmě, že tentokrát nebylo užito křížení, proto je ve výsledcích
méně kombinací.

Algoritmus VEGA dospěl k průměrnému jedinci s
bent skóre 8.1,
$-log_2{DP_{\max}} = 0.1$ tj. $LP_{\max} = 0.93303299153680741$,
$-log_2{DP_{\max}} = 0.5$ tj. $DP_{\max} = \frac{\sqrt{2}}{2}$,
bijektivní skóre 10.3,
SAC 0.0,
faktor větvení 1.0
a stupeň polynomu 1.1.
V tabulce \ref{tabVegaParametersCgp}
jsou uvedeny výsledky pro jednotlivé pravděpodobnosti mutace.
V grafu na obrázku \ref{cgpVegaParameters} jsou porovnány počty nadprůměrných .
Nejlépe vyšla pravděpodobnost mutace $pMut = 0.9$.
Na této reprezentaci vyšlo lépe bijektivní skóre, avšak o to jsou horší ostatní parametry.

\insertimage{cgpVegaParameters}{Parametry cgp pro VEGA}

Algoritmus SPEA dospěl k průměrnému jedinci s
bent skóre 21.8,
$-log_2{LP_{\max}} = 0.3$ tj. $LP_{\max} = 0.81225239635623558$,
$-log_2{DP_{\max}} = 0.8$ tj. $DP_{\max} = 0.57434917749851744$,
bijektivní skóre 12.7,
SAC 0.3,
faktor větvení 1.4
a stupeň polynomu 2.1.
Algoritmus Spea bez křížení nedospěl k žádné nadprůměrné kombinaci.
Na této reprezentaci je o něco lepší bijektivní skóre, a o to jsou horší ostatní parametry.

\section{Hledání složitějších s-boxů}
Bylo provedeno i hledání s-boxů o rozměrech $6 \times 6$, $6 \times 4$, $8 \times 8$ algoritmy
VEGA a SPEA, jak na permutacích tak na reprezentacích CGP, i686.
Opět se potvrdilo, že pokud je bijektivita zaručena reprezentací, dosáhlo se vysokých hodnot kritérií, s tím, že
vysoce hodnocených bylo několik a každý jedinec byl lepší v něčem jiném (pareto optimum).
Při bijektivitě jako zvláštní kritérium byly hodnoty všech kritérií malé.
Již na 30 generacích a velikosti populace 30 a počtu běhů 4, kdy z jednotlivých běhů byla vytvořena
sada nedominovaných výsledků.
Pro všechny výsledky v této nedominované sadě platilo pro rozměry $6 \times 6$ při
reprezentaci permutace, že $-log_2{LP_{\max}} = 2.38529015$ nebo $2.83007503$ pro algoritmus
SPEA, u algoritmu VEGA mělo několik jedinců $2$. Stupeň polynomu byl $4$ nebo $5$ u obou
algoritmů. Oba nalezly různé jedince s $-log_2{DP_{\max}} = 3.41503739$.
\begin{center}
\begin{minipage}{0.6\textwidth}
U algoritmu SPEA měl tento jedinec chromozóm:
\begin{verbatim}
5 49 43 44 3 30 35 59 40 29 13 53 11 4 2 18
39 41 38 10 57 36 7 20 17 37 47 24 5 55 62 51
22 60 26 16 14 0 21 63 48 1 34 19 27 52 42 61
12 50 46 54 32 31 58 8 15 6 9 56 33 28 23 45
\end{verbatim}

U algoritmu VEGA měl tento jedinec chromozóm:
\begin{verbatim}
50 35 21 23 3 38 15 25 26 11 8 62 32 42 47 51
57 28 24 48 58 54 19 56 55 13 60 2 43 0 52 20
46 29 14 59 31 17 9 41 30 22 40 4 49 10 37 27
63 45 5 33 6 39 44 61 18 16 34 53 1 36 7 12
\end{verbatim}
\end{minipage}
\end{center}

Obdobné výsledky byly i u počtu generací 200 a velikosti populace 200.
U s-boxů $6 \times 4$ byly všechny $LP_{\max} = 1$.
Některé s-boxy $8 \times 8$ docela rychle splnily podmínky\footnote{které jsou uvedeny v \cite{twofish}}
pro výběr $LP_{\max} \leq  0.0625$ a
$DP_{\max} \leq  \frac{10}{256}$
při hledání pomocí algoritmu SPEA, počtu generací 30 a velikosti populace 30.
Na obr. \ref{good88Box} je uvedený jeden z nich, který má
bent skóre 250,
$-log_2{LP_{\max}} = 4.18621874$ tj. $LP_{\max} = 0.05493164$,
$-log_2{DP_{\max}} = 4.67807198$ tj. $DP_{\max} \dot{=} \frac{10}{256}$,
bijektivní skóre 255,
SAC 0,
faktor větvení 2
a stupeň polynomu 7.
Zobrazení nuly na nulu je spíše nežádoucí, avšak to se týká jen u tohohle konkrétního jedince.

Ostatní řešení, co měla $-log_2{DP_{\max}} = 4$, měla o něco málo vyšší bent skóre.
K obdobným výsledkům dojde i náhodné prohledávání, ale potřebuje více běhů.
Náhodné permutace větších velikostí totiž mají nízké pravděpodobnosti
$LP_{\max}$ a $DP_{\max}$, jak je uvedeno v \cite{bigSboxes}. A hodně náhodných permutací má
$-log_2{DP_{\max}} = 4.67807198$ a $-log_2{DP_{\max}} = 4$.
K jedinci s $-log_2{LP_{\max}} = 4.18621874$\footnote{a dalšími parametry výše}
však evoluce dospěje evoluce s méně běhy, jinak řečeno je nutné méně náhodně vygenerovaných
jedinců.

\pagebreak
\section{Dvoufázové hledání}
Kvůli zaručení bijektivity u reprezentací CGP a i686 se jako další
možnost\footnote{po bijektivním skóre jako kritérium a omezení bijektivity v objektivní
funkci}
řešení nabízí rozdělit evoluci na dvě fáze.
V první fázi se na reprezentaci permutace najde nejlepší řešení. To je pak
použito jako trénovací množina pro symbolickou regresi pomocí reprezentací
CGP či i686.
Prvotně nastává problém s výběrem
vhodného kandidáta, který postoupí do druhé fáze. Symbolická regrese bude
muset být provedena několikrát pro více jedinců. Dále se bude muset přijmout
fakt, že i symbolická regrese má svá omezení a problém se na ni jen přenáší.
Symbolická regrese byla provedena
na reprezentacích CGP a i686 o rozměrech $2 \times 2$, $4 \times 4$, $6 \times 6$, $8 \times 8$
na různých permutacích jako trénovací množiny.
Trénovací množina je úplná, tj. obsahuje všechny vstupy a výstupy.

\subsubsection{Parametry symbolické regrese}
Již při počtu generací 50 a velikosti populace 1000 se v reprezentaci jako acyklický graf
s 3 sloupci, 3 řádky a 3 maximálními mutacemi, nalezne $3 \times 3$ za nutnosti
několika běhů. U složitějších s-boxů již záleží na tvaru funkce, např. pro
$f(x) =  x$ tj. $[1, 2, \cdots, n - 1]$, je bez problému nalezeno řešení i pro
$8 \times 8$.

Vedlejším produktem tohoto rozdělení může být optimální realizace, avšak
základním kritériem musí být přesnost. Nejoptimálnější se vybere až z několika
přesných řešení.

\section{Celkové zhodnocení výsledku}
Při jednokriteriální optimalizaci na provedených experimentech
dochází jak náhodné prohledávání tak genetický algoritmus
k téměř obdobným výsledkům. Pouze v některých případech je výrazněji lepší buď jedno nebo druhé.
V multikriteriálním prohledávání pomocí algoritmů VEGA a SPEA
lze snadněji zajistit lepší hodnoty více kritérií souběžně
než u náhodného prohledávání. Zatímco náhodné prohledávání je algoritmicky méně složitější,
genetický algoritmus nepotřebuje tolik běhů (náhodných inicializací jedinců na počátku).
Tedy jsou případy reprezentací jako permutace, kdy lze evoluční algoritmy použít.
Zaručit bijektivitu lze spolehlivě jen u permutací. U ostatních reprezentací se při multikriteriálním prohledávání
po přidání bijektivního skóre výrazně zhoršují hodnoty ostatních kritérií. A pokud
pro všechny nebijektivní jedince bude hodnota každého kritéria nulová, či-li napřed je
proveden test na bijektivitu a poté je teprve proveden výpočet kritéria, poskytuje
ještě horší výsledky. Nejlépe by bylo, kdyby pro reprezentace CGP a i686 existovaly
genetické operátory, které bijektivitu zaručí.
U těchto reprezentací se projevuje složitost jejich struktur (mnoho možností
pro mutace a křížení a častost opakování toho samého jedince), a tím pádem
evoluční algoritmus hledá řešení obtížněji \cite{prednaskyBIN}.
Obdobné závěry lze konstatovat i u hledání minimální realizace (podle počtu použitých hradel) u reprezentace CGP, kdy
je populace zanesena minimálními nefunkčními obvody.
Rozdělení na dvě fáze je pouhé přenesení na problém symbolické regrese. Minimální
realizace musí být vybrána, až po splnění kritérií bezpečnosti.
V nalezených zdrojích se nezabývají multikriteriálním prohledáváním s více kritérii
bezpečnosti.
V \cite{paretoHardware} se zabývají optimální strukturou hardware, avšak hodnoty kritérií
bezpečnosti nejsou zmíněny.
